% ============================================================
% Módulos Funcionales del Sistema — Cap. 4.1.2
% 7 Módulos lógicos — Paquetes UML simplificado (solo lógica)
% Proyecto BTH V2.0 — Sistema Web de Inscripción MTP San Julián
% Formato: TikZ vectorial editable (APA 7)
% Uso: % ============================================================
% Módulos Funcionales del Sistema — Cap. 4.1.2
% 7 Módulos lógicos — Paquetes UML simplificado (solo lógica)
% Proyecto BTH V2.0 — Sistema Web de Inscripción MTP San Julián
% Formato: TikZ vectorial editable (APA 7)
% Uso: % ============================================================
% Módulos Funcionales del Sistema — Cap. 4.1.2
% 7 Módulos lógicos — Paquetes UML simplificado (solo lógica)
% Proyecto BTH V2.0 — Sistema Web de Inscripción MTP San Julián
% Formato: TikZ vectorial editable (APA 7)
% Uso: % ============================================================
% Módulos Funcionales del Sistema — Cap. 4.1.2
% 7 Módulos lógicos — Paquetes UML simplificado (solo lógica)
% Proyecto BTH V2.0 — Sistema Web de Inscripción MTP San Julián
% Formato: TikZ vectorial editable (APA 7)
% Uso: \input{imagenes/modulos_paquetes} dentro de figure
% ============================================================
\begin{tikzpicture}[
    modulo/.style={
        rectangle, rounded corners=3.5pt,
        draw=primary!75, fill=white,
        text=darktext, align=center,
        font=\scriptsize,
        text width=3.42cm, minimum height=1.35cm,
        inner sep=4pt
    },
    moduloDestacado/.style={
        rectangle, rounded corners=3.5pt,
        draw=primary, fill=primary,
        text=white, align=center,
        font=\scriptsize\bfseries,
        text width=3.42cm, minimum height=1.35cm,
        inner sep=4pt
    },
    moduloLight/.style={
        rectangle, rounded corners=3.5pt,
        draw=primary!55, fill=secondary!11,
        text=darktext, align=center,
        font=\scriptsize,
        text width=3.42cm, minimum height=1.35cm,
        inner sep=4pt
    },
    flecha/.style={
        -{Stealth[length=2.6mm,width=2mm]},
        thick, draw=primary!80
    },
    flechaDashed/.style={
        -{Stealth[length=2.6mm,width=2mm]},
        thick, draw=primary!55, dashed
    },
    etiqueta/.style={
        font=\scriptsize\bfseries, text=primary, align=center
    },
    sistema/.style={
        rectangle, rounded corners=7pt,
        draw=primary!35, thick,
        inner sep=6pt
    }
]

% --- Contenedor Sistema ---
\node[sistema, minimum width=13.2cm, minimum height=6.35cm] (sistema) at (0, 0.15) {};

% Título sistema
\node[etiqueta, fill=white, inner sep=2pt] at (0, 3.55) {SISTEMA WEB BTH V2.0 — MÓDULOS FUNCIONALES};

% --- NIVEL SUPERIOR: Preinscripción (entrada) ---
\node[moduloDestacado] (M1) at (0, 2.45) {\textbf{M1. Formulario}\\Dinámico y\\Preinscripción\\\footnotesize\normalfont\textit{7 especialidades}};

% --- NIVEL MEDIO: Núcleo transaccional + Recovery ---
\node[modulo] (M2) at (-4.15, 0.70) {\textbf{M2. Control}\\Transaccional\\de Vacantes\\\footnotesize\normalfont\textit{ACID row-level}};
\node[moduloLight] (M3) at (0, 0.70) {\textbf{M3. Recuperación}\\Autónoma de\\Credenciales\\\footnotesize\normalfont\textit{C.I.+fecha <$2$ min}};
\node[moduloLight] (M6) at (4.15, 0.70) {\textbf{M6. Generador}\\Comprobantes\\PDF\\\footnotesize\normalfont\textit{código verificación}};

% --- NIVEL INFERIOR: Gestión académica ---
\node[modulo] (M4) at (-4.15, -1.25) {\textbf{M4. Matriculación}\\Documental\\\footnotesize\normalfont\textit{libretas 1ro-3ro,}\\[-2pt]\footnotesize\normalfont\textit{C.I. tutor, carnet, cert.}};
\node[modulo] (M5) at (0, -1.25) {\textbf{M5. Asignación}\\Cursos y Paralelos\\\footnotesize\normalfont\textit{por U.E. procedencia}\\[-2pt]\footnotesize\normalfont\textit{4to A / 4to B}};
\node[modulo] (M7) at (4.15, -1.25) {\textbf{M7. Panel Control}\\y Reportes\\\footnotesize\normalfont\textit{vacantes, \% matrícula}};

% --- Conexiones principales (flujo) ---
% M1 -> M2 (reserva)
\draw[flecha] (M1.south) -- (-2.05, 1.42) -- (M2.north);
\node[font=\tiny, text=primary!80, fill=white, inner sep=1pt] at (-2.35, 1.65) {reserva};

% M1 -> M6 (comprobante)
\draw[flechaDashed] (M1.east) -- (4.85, 1.42) -- (M6.north);
\node[font=\tiny, text=primary!65, fill=white, inner sep=1pt] at (3.15, 1.90) {PDF};

% M1 -> M3 (soporte, recuperación)
\draw[flechaDashed] (M1.south) -- (0, 1.42) -- (M3.north);

% M2 -> M4 (matriculación tras reserva)
\draw[flecha] (M2.south) -- (M4.north);

% M2 -> M5 (asignación tras matrícula)
\draw[flecha] (M2.south) -- (-1.55, -0.20) -- (M5.north);
% M4 -> M5
\draw[flecha] (M4.east) -- (M5.west);
\node[font=\tiny, text=primary!80, fill=white, inner sep=1pt] at (-2.05, -0.85) {exp. ok};

% M5 -> M7 y M6 -> M7 (reportes)
\draw[flechaDashed] (M5.east) -- (M7.west);
\draw[flechaDashed] (M6.south) -- (4.15, -0.05) -- (M7.north);

% M2 -> M7 (vacantes tiempo real)
\draw[flechaDashed] (M2.east) -- (2.10, -0.20) -- (4.15, -0.20) -- (M7.north west);

% --- Leyenda inferior ---
\node[font=\tiny\itshape, text=darktext!65, align=center] at (0, -2.85) {Flujo 5 fases Cap.~4.1.3: (1) Preinscripción → (2) Recovery → (3) Comprobante → (4) Matriculación → (5) Asignación \textbullet{} M7 transversal en tiempo real};
\draw[primary!80, thick] ( -6.0, -2.55) -- (6.0, -2.55);

% --- Indicadores pequeños ---
\node[font=\tiny, text=primary!50, align=left] at (-6.15, 3.15) {\textbf{—} flujo};
\node[font=\tiny, text=primary!50, align=left] at (-6.15, 2.85) {\textbf{- -} consulta};

\end{tikzpicture}
 dentro de figure
% ============================================================
\begin{tikzpicture}[
    modulo/.style={
        rectangle, rounded corners=3.5pt,
        draw=primary!75, fill=white,
        text=darktext, align=center,
        font=\scriptsize,
        text width=3.42cm, minimum height=1.35cm,
        inner sep=4pt
    },
    moduloDestacado/.style={
        rectangle, rounded corners=3.5pt,
        draw=primary, fill=primary,
        text=white, align=center,
        font=\scriptsize\bfseries,
        text width=3.42cm, minimum height=1.35cm,
        inner sep=4pt
    },
    moduloLight/.style={
        rectangle, rounded corners=3.5pt,
        draw=primary!55, fill=secondary!11,
        text=darktext, align=center,
        font=\scriptsize,
        text width=3.42cm, minimum height=1.35cm,
        inner sep=4pt
    },
    flecha/.style={
        -{Stealth[length=2.6mm,width=2mm]},
        thick, draw=primary!80
    },
    flechaDashed/.style={
        -{Stealth[length=2.6mm,width=2mm]},
        thick, draw=primary!55, dashed
    },
    etiqueta/.style={
        font=\scriptsize\bfseries, text=primary, align=center
    },
    sistema/.style={
        rectangle, rounded corners=7pt,
        draw=primary!35, thick,
        inner sep=6pt
    }
]

% --- Contenedor Sistema ---
\node[sistema, minimum width=13.2cm, minimum height=6.35cm] (sistema) at (0, 0.15) {};

% Título sistema
\node[etiqueta, fill=white, inner sep=2pt] at (0, 3.55) {SISTEMA WEB BTH V2.0 — MÓDULOS FUNCIONALES};

% --- NIVEL SUPERIOR: Preinscripción (entrada) ---
\node[moduloDestacado] (M1) at (0, 2.45) {\textbf{M1. Formulario}\\Dinámico y\\Preinscripción\\\footnotesize\normalfont\textit{7 especialidades}};

% --- NIVEL MEDIO: Núcleo transaccional + Recovery ---
\node[modulo] (M2) at (-4.15, 0.70) {\textbf{M2. Control}\\Transaccional\\de Vacantes\\\footnotesize\normalfont\textit{ACID row-level}};
\node[moduloLight] (M3) at (0, 0.70) {\textbf{M3. Recuperación}\\Autónoma de\\Credenciales\\\footnotesize\normalfont\textit{C.I.+fecha <$2$ min}};
\node[moduloLight] (M6) at (4.15, 0.70) {\textbf{M6. Generador}\\Comprobantes\\PDF\\\footnotesize\normalfont\textit{código verificación}};

% --- NIVEL INFERIOR: Gestión académica ---
\node[modulo] (M4) at (-4.15, -1.25) {\textbf{M4. Matriculación}\\Documental\\\footnotesize\normalfont\textit{libretas 1ro-3ro,}\\[-2pt]\footnotesize\normalfont\textit{C.I. tutor, carnet, cert.}};
\node[modulo] (M5) at (0, -1.25) {\textbf{M5. Asignación}\\Cursos y Paralelos\\\footnotesize\normalfont\textit{por U.E. procedencia}\\[-2pt]\footnotesize\normalfont\textit{4to A / 4to B}};
\node[modulo] (M7) at (4.15, -1.25) {\textbf{M7. Panel Control}\\y Reportes\\\footnotesize\normalfont\textit{vacantes, \% matrícula}};

% --- Conexiones principales (flujo) ---
% M1 -> M2 (reserva)
\draw[flecha] (M1.south) -- (-2.05, 1.42) -- (M2.north);
\node[font=\tiny, text=primary!80, fill=white, inner sep=1pt] at (-2.35, 1.65) {reserva};

% M1 -> M6 (comprobante)
\draw[flechaDashed] (M1.east) -- (4.85, 1.42) -- (M6.north);
\node[font=\tiny, text=primary!65, fill=white, inner sep=1pt] at (3.15, 1.90) {PDF};

% M1 -> M3 (soporte, recuperación)
\draw[flechaDashed] (M1.south) -- (0, 1.42) -- (M3.north);

% M2 -> M4 (matriculación tras reserva)
\draw[flecha] (M2.south) -- (M4.north);

% M2 -> M5 (asignación tras matrícula)
\draw[flecha] (M2.south) -- (-1.55, -0.20) -- (M5.north);
% M4 -> M5
\draw[flecha] (M4.east) -- (M5.west);
\node[font=\tiny, text=primary!80, fill=white, inner sep=1pt] at (-2.05, -0.85) {exp. ok};

% M5 -> M7 y M6 -> M7 (reportes)
\draw[flechaDashed] (M5.east) -- (M7.west);
\draw[flechaDashed] (M6.south) -- (4.15, -0.05) -- (M7.north);

% M2 -> M7 (vacantes tiempo real)
\draw[flechaDashed] (M2.east) -- (2.10, -0.20) -- (4.15, -0.20) -- (M7.north west);

% --- Leyenda inferior ---
\node[font=\tiny\itshape, text=darktext!65, align=center] at (0, -2.85) {Flujo 5 fases Cap.~4.1.3: (1) Preinscripción → (2) Recovery → (3) Comprobante → (4) Matriculación → (5) Asignación \textbullet{} M7 transversal en tiempo real};
\draw[primary!80, thick] ( -6.0, -2.55) -- (6.0, -2.55);

% --- Indicadores pequeños ---
\node[font=\tiny, text=primary!50, align=left] at (-6.15, 3.15) {\textbf{—} flujo};
\node[font=\tiny, text=primary!50, align=left] at (-6.15, 2.85) {\textbf{- -} consulta};

\end{tikzpicture}
 dentro de figure
% ============================================================
\begin{tikzpicture}[
    modulo/.style={
        rectangle, rounded corners=3.5pt,
        draw=primary!75, fill=white,
        text=darktext, align=center,
        font=\scriptsize,
        text width=3.42cm, minimum height=1.35cm,
        inner sep=4pt
    },
    moduloDestacado/.style={
        rectangle, rounded corners=3.5pt,
        draw=primary, fill=primary,
        text=white, align=center,
        font=\scriptsize\bfseries,
        text width=3.42cm, minimum height=1.35cm,
        inner sep=4pt
    },
    moduloLight/.style={
        rectangle, rounded corners=3.5pt,
        draw=primary!55, fill=secondary!11,
        text=darktext, align=center,
        font=\scriptsize,
        text width=3.42cm, minimum height=1.35cm,
        inner sep=4pt
    },
    flecha/.style={
        -{Stealth[length=2.6mm,width=2mm]},
        thick, draw=primary!80
    },
    flechaDashed/.style={
        -{Stealth[length=2.6mm,width=2mm]},
        thick, draw=primary!55, dashed
    },
    etiqueta/.style={
        font=\scriptsize\bfseries, text=primary, align=center
    },
    sistema/.style={
        rectangle, rounded corners=7pt,
        draw=primary!35, thick,
        inner sep=6pt
    }
]

% --- Contenedor Sistema ---
\node[sistema, minimum width=13.2cm, minimum height=6.35cm] (sistema) at (0, 0.15) {};

% Título sistema
\node[etiqueta, fill=white, inner sep=2pt] at (0, 3.55) {SISTEMA WEB BTH V2.0 — MÓDULOS FUNCIONALES};

% --- NIVEL SUPERIOR: Preinscripción (entrada) ---
\node[moduloDestacado] (M1) at (0, 2.45) {\textbf{M1. Formulario}\\Dinámico y\\Preinscripción\\\footnotesize\normalfont\textit{7 especialidades}};

% --- NIVEL MEDIO: Núcleo transaccional + Recovery ---
\node[modulo] (M2) at (-4.15, 0.70) {\textbf{M2. Control}\\Transaccional\\de Vacantes\\\footnotesize\normalfont\textit{ACID row-level}};
\node[moduloLight] (M3) at (0, 0.70) {\textbf{M3. Recuperación}\\Autónoma de\\Credenciales\\\footnotesize\normalfont\textit{C.I.+fecha <$2$ min}};
\node[moduloLight] (M6) at (4.15, 0.70) {\textbf{M6. Generador}\\Comprobantes\\PDF\\\footnotesize\normalfont\textit{código verificación}};

% --- NIVEL INFERIOR: Gestión académica ---
\node[modulo] (M4) at (-4.15, -1.25) {\textbf{M4. Matriculación}\\Documental\\\footnotesize\normalfont\textit{libretas 1ro-3ro,}\\[-2pt]\footnotesize\normalfont\textit{C.I. tutor, carnet, cert.}};
\node[modulo] (M5) at (0, -1.25) {\textbf{M5. Asignación}\\Cursos y Paralelos\\\footnotesize\normalfont\textit{por U.E. procedencia}\\[-2pt]\footnotesize\normalfont\textit{4to A / 4to B}};
\node[modulo] (M7) at (4.15, -1.25) {\textbf{M7. Panel Control}\\y Reportes\\\footnotesize\normalfont\textit{vacantes, \% matrícula}};

% --- Conexiones principales (flujo) ---
% M1 -> M2 (reserva)
\draw[flecha] (M1.south) -- (-2.05, 1.42) -- (M2.north);
\node[font=\tiny, text=primary!80, fill=white, inner sep=1pt] at (-2.35, 1.65) {reserva};

% M1 -> M6 (comprobante)
\draw[flechaDashed] (M1.east) -- (4.85, 1.42) -- (M6.north);
\node[font=\tiny, text=primary!65, fill=white, inner sep=1pt] at (3.15, 1.90) {PDF};

% M1 -> M3 (soporte, recuperación)
\draw[flechaDashed] (M1.south) -- (0, 1.42) -- (M3.north);

% M2 -> M4 (matriculación tras reserva)
\draw[flecha] (M2.south) -- (M4.north);

% M2 -> M5 (asignación tras matrícula)
\draw[flecha] (M2.south) -- (-1.55, -0.20) -- (M5.north);
% M4 -> M5
\draw[flecha] (M4.east) -- (M5.west);
\node[font=\tiny, text=primary!80, fill=white, inner sep=1pt] at (-2.05, -0.85) {exp. ok};

% M5 -> M7 y M6 -> M7 (reportes)
\draw[flechaDashed] (M5.east) -- (M7.west);
\draw[flechaDashed] (M6.south) -- (4.15, -0.05) -- (M7.north);

% M2 -> M7 (vacantes tiempo real)
\draw[flechaDashed] (M2.east) -- (2.10, -0.20) -- (4.15, -0.20) -- (M7.north west);

% --- Leyenda inferior ---
\node[font=\tiny\itshape, text=darktext!65, align=center] at (0, -2.85) {Flujo 5 fases Cap.~4.1.3: (1) Preinscripción → (2) Recovery → (3) Comprobante → (4) Matriculación → (5) Asignación \textbullet{} M7 transversal en tiempo real};
\draw[primary!80, thick] ( -6.0, -2.55) -- (6.0, -2.55);

% --- Indicadores pequeños ---
\node[font=\tiny, text=primary!50, align=left] at (-6.15, 3.15) {\textbf{—} flujo};
\node[font=\tiny, text=primary!50, align=left] at (-6.15, 2.85) {\textbf{- -} consulta};

\end{tikzpicture}
 dentro de figure
% ============================================================
\begin{tikzpicture}[
    modulo/.style={
        rectangle, rounded corners=3.5pt,
        draw=primary!75, fill=white,
        text=darktext, align=center,
        font=\scriptsize,
        text width=3.42cm, minimum height=1.35cm,
        inner sep=4pt
    },
    moduloDestacado/.style={
        rectangle, rounded corners=3.5pt,
        draw=primary, fill=primary,
        text=white, align=center,
        font=\scriptsize\bfseries,
        text width=3.42cm, minimum height=1.35cm,
        inner sep=4pt
    },
    moduloLight/.style={
        rectangle, rounded corners=3.5pt,
        draw=primary!55, fill=secondary!11,
        text=darktext, align=center,
        font=\scriptsize,
        text width=3.42cm, minimum height=1.35cm,
        inner sep=4pt
    },
    flecha/.style={
        -{Stealth[length=2.6mm,width=2mm]},
        thick, draw=primary!80
    },
    flechaDashed/.style={
        -{Stealth[length=2.6mm,width=2mm]},
        thick, draw=primary!55, dashed
    },
    etiqueta/.style={
        font=\scriptsize\bfseries, text=primary, align=center
    },
    sistema/.style={
        rectangle, rounded corners=7pt,
        draw=primary!35, thick,
        inner sep=6pt
    }
]

% --- Contenedor Sistema ---
\node[sistema, minimum width=13.2cm, minimum height=6.35cm] (sistema) at (0, 0.15) {};

% Título sistema
\node[etiqueta, fill=white, inner sep=2pt] at (0, 3.55) {SISTEMA WEB BTH V2.0 — MÓDULOS FUNCIONALES};

% --- NIVEL SUPERIOR: Preinscripción (entrada) ---
\node[moduloDestacado] (M1) at (0, 2.45) {\textbf{M1. Formulario}\\Dinámico y\\Preinscripción\\\footnotesize\normalfont\textit{7 especialidades}};

% --- NIVEL MEDIO: Núcleo transaccional + Recovery ---
\node[modulo] (M2) at (-4.15, 0.70) {\textbf{M2. Control}\\Transaccional\\de Vacantes\\\footnotesize\normalfont\textit{ACID row-level}};
\node[moduloLight] (M3) at (0, 0.70) {\textbf{M3. Recuperación}\\Autónoma de\\Credenciales\\\footnotesize\normalfont\textit{C.I.+fecha <$2$ min}};
\node[moduloLight] (M6) at (4.15, 0.70) {\textbf{M6. Generador}\\Comprobantes\\PDF\\\footnotesize\normalfont\textit{código verificación}};

% --- NIVEL INFERIOR: Gestión académica ---
\node[modulo] (M4) at (-4.15, -1.25) {\textbf{M4. Matriculación}\\Documental\\\footnotesize\normalfont\textit{libretas 1ro-3ro,}\\[-2pt]\footnotesize\normalfont\textit{C.I. tutor, carnet, cert.}};
\node[modulo] (M5) at (0, -1.25) {\textbf{M5. Asignación}\\Cursos y Paralelos\\\footnotesize\normalfont\textit{por U.E. procedencia}\\[-2pt]\footnotesize\normalfont\textit{4to A / 4to B}};
\node[modulo] (M7) at (4.15, -1.25) {\textbf{M7. Panel Control}\\y Reportes\\\footnotesize\normalfont\textit{vacantes, \% matrícula}};

% --- Conexiones principales (flujo) ---
% M1 -> M2 (reserva)
\draw[flecha] (M1.south) -- (-2.05, 1.42) -- (M2.north);
\node[font=\tiny, text=primary!80, fill=white, inner sep=1pt] at (-2.35, 1.65) {reserva};

% M1 -> M6 (comprobante)
\draw[flechaDashed] (M1.east) -- (4.85, 1.42) -- (M6.north);
\node[font=\tiny, text=primary!65, fill=white, inner sep=1pt] at (3.15, 1.90) {PDF};

% M1 -> M3 (soporte, recuperación)
\draw[flechaDashed] (M1.south) -- (0, 1.42) -- (M3.north);

% M2 -> M4 (matriculación tras reserva)
\draw[flecha] (M2.south) -- (M4.north);

% M2 -> M5 (asignación tras matrícula)
\draw[flecha] (M2.south) -- (-1.55, -0.20) -- (M5.north);
% M4 -> M5
\draw[flecha] (M4.east) -- (M5.west);
\node[font=\tiny, text=primary!80, fill=white, inner sep=1pt] at (-2.05, -0.85) {exp. ok};

% M5 -> M7 y M6 -> M7 (reportes)
\draw[flechaDashed] (M5.east) -- (M7.west);
\draw[flechaDashed] (M6.south) -- (4.15, -0.05) -- (M7.north);

% M2 -> M7 (vacantes tiempo real)
\draw[flechaDashed] (M2.east) -- (2.10, -0.20) -- (4.15, -0.20) -- (M7.north west);

% --- Leyenda inferior ---
\node[font=\tiny\itshape, text=darktext!65, align=center] at (0, -2.85) {Flujo 5 fases Cap.~4.1.3: (1) Preinscripción → (2) Recovery → (3) Comprobante → (4) Matriculación → (5) Asignación \textbullet{} M7 transversal en tiempo real};
\draw[primary!80, thick] ( -6.0, -2.55) -- (6.0, -2.55);

% --- Indicadores pequeños ---
\node[font=\tiny, text=primary!50, align=left] at (-6.15, 3.15) {\textbf{—} flujo};
\node[font=\tiny, text=primary!50, align=left] at (-6.15, 2.85) {\textbf{- -} consulta};

\end{tikzpicture}
